% !TeX root = ../navier-stokes-manuscript.tex
\subsection{The equation and the target}\label{sub:TheEquationAndTheTarget}

Fefferman's formulation of the whole-space problem begins with an initial
velocity that leaves no roughness to blame later \parencite{Fefferman2000}.
It is smooth, divergence-free, and every one of its spatial derivatives dies
faster than any polynomial as \(|x|\) moves out through the fluid.  Write this
class as
\begin{equation}\label{eq:FeffermanDatumClass}
  \Ssigma
  :=
  \left\{
  a\in C^\infty(\R^3;\R^3):
  \begin{array}{l}
    \nabla\cdot a=0,\\
    \forall\alpha\in\Nzero^3\;
    \forall K\in\Nzero\;
    \exists C_{\alpha,K}<\infty\;
    \forall x\in\R^3,\\
    |\partial_x^\alpha a(x)|
    \leq C_{\alpha,K}(1+|x|)^{-K}
  \end{array}
  \right\}.
\end{equation}
Read the condition in the order it is written.  Choose any spatial derivative
\(\alpha\).  Choose any algebraic decay rate \(K\), however large.  The datum
still has a finite constant \(C_{\alpha,K}\) that controls that derivative
everywhere in space.  The constant changes when either choice changes; it does
not change with \(x\).  This is the divergence-free Schwartz class.

For a finite spatial height \(m\), let
\[
  H^m_\sigma(\R^3)
  :=
  \{v\in H^m(\R^3;\R^3):\nabla\cdot v=0\}.
\]
The decay in \eqref{eq:FeffermanDatumClass} places the same datum in every one
of these spaces.  To see it, fix \(K>3/2\).  For \(|\alpha|\leq m\),
\[
  \norm{\partial_x^\alpha a}_2^2
  \leq
  C_{\alpha,K}^2
  \int_{\R^3}(1+|x|)^{-2K}\,dx
  <\infty.
\]
Only finitely many derivatives occur below a chosen finite \(m\), and hence
\begin{equation}\label{eq:SchwartzToEveryFiniteSobolevHeight}
  \Ssigma\subset H^m_\sigma(\R^3)
  \qquad\text{for every finite }m.
\end{equation}
At time zero, then, every finite Sobolev height is already present.  The
regularity problem asks whether the evolution can destroy that complete
smooth structure in finite time.

Fix one \(u_0\in\Ssigma\) and one viscosity \(\nu>0\).  The velocity \(u\) and
the pressure \(p\) obey one coupled law:
\begin{equation}\label{eq:NavierStokesSystem}
  \partial_tu+(u\cdot\nabla)u+\nabla p=\nu\Delta u,
  \qquad \nabla\cdot u=0,
  \qquad u(\cdot,0)=u_0.
\end{equation}
The four parts should be kept together.  The time derivative records the
response of the field.  Transport lets the velocity carry itself through the
velocity it has already created.  Incompressibility forces pressure to answer
for that motion across the surrounding field.  Viscosity acts on the spatial
curvature of the same velocity.  There is one fluid here, and every later
derivative row is another reading of this law.

Local theory starts that fluid smoothly.  Uniqueness then joins the local
pieces into one maximal classical history on an interval \([0,T_*)\), where
\(T_*\in(0,\infty]\).  The Millennium Prize problem is the claim that this
interval never has a finite right edge.  In the notation above, its target is
\begin{equation}\label{eq:WholeSpaceRegularityTarget}
  T_*=\infty,
  \qquad
  u,p\in C^\infty\!\bigl(\R^3\times[0,\infty)\bigr),
  \qquad
  \sup_{t\geq0}\int_{\R^3}|u(x,t)|^2\,dx<\infty.
\end{equation}
Smoothness includes the last slice of every finite interval.  A field can be
smooth on every shorter interval \([0,T]\), \(T<T_*\), while one of its
derivatives grows without a finite trace as \(t\uparrow T_*\).  Interior
smoothness and arrival at the endpoint are different statements.  The proof
has to carry the derivatives to that face or show that the face cannot exist.
